\section{Top project works}
\begin{enumerate}
	\item System identification of an airplane using Physics Informed Neural Networks (PINNs)
		\begin{itemize}
			\item System identification for the longitudinal dynamics of Hansa airplane was performed in the work.
			\item A linear aerodynamic model and the dynamics governing equations in wind axis reference frame were implemented.
			\item Python programming language with Tensorflow library was used in the work. Estimations were performed using both synthetic and experimental data.
			\item A new PINN architecture, System-Identification PINNs (SI-PINNs) was developed,
				overcoming limitations faced by existing models in the literature.
			\item Performance of SI-PINNs was compared with classical methods, and the advantages of SI-PINNs in handling sparse and faulty data were demonstrated.
			\item SI-PINN is proposed as an alternative to classical methods in the aircraft system identification, and the experimentations can be geared
				towards sparse but accurate data instead of huge and faulty.
		\end{itemize}

	\item Development of custom OpenFOAM solver for the wave equation
		\begin{itemize}
			\item A custom application named \textit{waveFoam} was developed using OpenFOAM C++ to
				solve the transient wave equation.
			\item The solver was developed to handle a maximum of upto 4 dimensional solution field (3 spatial and 1 temporal) and embodied with all features that a basic OpenFOAM application has.
			\item A custom foam dictionary was made to feed the wave speed property needed by the solver that can access it during execution.
			\item A new field file \textit{h} was created that records the amplitude of waves in the domain.
			\item The solution files were saved in FOAM format to post-process directly using ParaView.
			\item The code was submitted as a tutorial for the blooming programmers in OpenFOAM community. \url{https://github.com/UnnamedMoose/BasicOpenFOAMProgrammingTutorials}
		\end{itemize}

	\item Development of incompressible single-phase massless particle tracking application in OpenFOAM
		\begin{itemize}
			\item There are few particle tracking applications in OpenFOAM v1806 but none of them are massless,
				hence this application was developed.
			\item The application tracks particles by displacing them with integration of fluid velocity
				over a given timestep. The particles are represented as points in 3D domain.
			\item A new foam dictionary, \textit{particleTracksDict} was developed that stores the
				initial position of particles, which are defined using either a list of coordinates or
				by defining the name of a boundary.
			\item Back or forward tracing will be determined by the application based on the specified
				inlet or outlet boundary.
			\item The application writes residence time and distance travelled by each particle to a
				post-processing file and the path of each particle is exported as vtk file for ease
				of visualization in ParaView.
			\item The work was submitted as a tutorial to the OpenFOAM community.\url{https://github.com/UnnamedMoose/BasicOpenFOAMProgrammingTutorials}
		\end{itemize}
	\item Numerical analysis of a laminar incompressible viscous flow over flat plate and the driven cavity flow using
		Lattice Boltzmann Method (LBM)
		\begin{itemize}
			\item Flow simulation was carried using air at standard condition and for Reynolds numbers of 10 and 1000
				for the flat plate and cavity flow, respectively.
			\item The Lattice Boltzmann Method (LBM), a mesoscopic approach was implemented in the work.
			\item The code was developed on MATLAB and Python languages.
			\item Validation was performed using Blasius solution for flat plate and the benchmark results from the literature for the cavity flow.
			\item The performance of LBM was compared to the finite difference method (FDM) for both the problems.
		\end{itemize}

	\item Development of a two-phase fluid flow solver for rectangular structured domains
		\begin{itemize}
			\item The solver was developed using staggered-finite volume method using FORTRAN-2008 for
				computation, NetCDF for data handling and Python 3.6 for post-processing.
			\item Immersed interface model and the Eulerian Volume of Fluid (VOF) method were
				implemented for the two phase flow in the work.
			\item Continuum surface force model was implemented to include surface tension forces in the computation.
			\item Pressure correction technique with modified SIMPLE algorithm was used to
				solve the incompressible governing equations.
			\item Drop at 0-gravity and the Rayleigh-Taylor instability problems were solved for the validation of the solver.
		\end{itemize}

	\item Numerical simulation of Poiseuille flow
		\begin{itemize}
			\item Axisymmetric Navier-Stokes equations using finite volume method using vorticity-streamfunction formulation were used in the work.
			\item CFD code was developed using MATLAB and Python. A simple GUI was developed using Python language.
			\item Validation was performed using Hagen-Poiseuille equation and comparing velocity profile with
				theoretical variation.
		\end{itemize}
\end{enumerate}

%------------------------------------------------------------------------------
\section{Github page}
This page contais some of my code works
\url{https://github.com/Ramkumar47}

%------------------------------------------------------------------------------
\section{References}
\begin{itemize}
	\item Dr. Madhukar M. Rao, Technical director, ACRi Info Tech, Bangalore. \href{mailto: madhukar.rao@acricfd.com}{madhukar.rao@acricfd.com}
	\item Dr. Devendra Prakash Ghate, Assistant Professor, Indian Institute of Space Science and Technology. \href{mailto: devendra.ghate@iist.ac.in}{devendra.ghate@iist.ac.in}
\end{itemize}

%------------------------------------------------------------------------------
\section{Declaration}
I, hereby, declare that the above given information is true to the best of my knowledge.
\vspace{2cm}



\begin{minipage}[t]{0.7\linewidth}
	Date:
\end{minipage}
\begin{minipage}[t]{0.5\linewidth}
	signature:
	(Ramkumar S)
\end{minipage}
